%% ma: L12-B08-0146
%% lop: 12
%% chuong: 2
%% bai: 8
%% dang: 12.8.49
%% loai: TN4
%% muc: VD
%% dap_an: "C"
%% nguon: "CĐ2.3.pdf giáo viên gửi 27/10/2026 (M07, Lê Minh Khôi – Nguyễn Anh Khôi, Chương II Chủ đề 3 Biểu thức tọa độ của các phép toán vectơ), Câu 62"
%% kien_thuc: "Bài 8 (điều kiện vuông góc, biểu thức tọa độ), lớp 10 (hệ phương trình bậc nhất)"
%% trang_thai: cho-duyet
%% ly_do: "Dạng mới 12.8.49 đề xuất, chờ giáo viên duyệt"
\begin{ex}
Trong không gian $Oxyz$, cho tam giác $IDF$ có tọa độ các đỉnh là $I(1;3;4)$, $D(3;6;7)$ và $F(-5;-3;-1)$. Tọa độ trực tâm $H$ của tam giác $IDF$ là
\choice{$H\left(\dfrac{5419}{109};\dfrac{753}{109};\dfrac{3659}{109}\right)$}{$H\left(-\dfrac{5419}{109};\dfrac{753}{109};\dfrac{3659}{109}\right)$}{\True $H\left(\dfrac{5419}{109};-\dfrac{753}{109};-\dfrac{3659}{109}\right)$}{$H\left(\dfrac{5419}{109};-\dfrac{753}{109};\dfrac{3659}{109}\right)$}
\loigiai{$\overrightarrow{IH}=s\overrightarrow{ID}+t\overrightarrow{IF}$ với $\overrightarrow{ID}=(2;3;3)$, $\overrightarrow{IF}=(-6;-6;-5)$; giải hệ $\overrightarrow{IH}\cdot\overrightarrow{DF}=0$, $\overrightarrow{DH}\cdot\overrightarrow{IF}=0$.

$H\left(\dfrac{5419}{109};-\dfrac{753}{109};-\dfrac{3659}{109}\right)$.}
\end{ex}
